%! TEX root = main.tex

\subsection{연구 배경}

데이터는 산업과 학계에서 원자재와 같은 중요한 전략 자산이 되었으며, 특히 생성형 AI 서비스의 대중화를 계기로 고품질 데이터에 대한 수요가 폭증하고 있다. 그중에서도 의료 분야에서 생산되는 데이터는 개인의 생명 및 건강과 직결된 가장 민감한 정보로, 그 활용 가치만큼이나 매우 강력한 수준의 보안이 요구된다. 그러나 의료 데이터가 지닌 높은 가치는 곧바로 공격자들의 표적이 되고 있다. 최신 2025년 데이터 침해 비용 보고서에 따르면, 의료 산업은 14년 연속으로 전 세계에서 데이터 유출 피해 비용이 가장 높은 산업으로 기록되었으며, 침해로 인한 평균 비용은 무려 742만 달러(USD)에 달한다\cite{ibm2025cost}. 공격자들은 신분 도용이나 보험 사기 등의 금융 범죄를 목적으로 환자의 개인식별정보(PII)와 의료 기록을 집중적으로 노리고 있다. 더욱 치명적인 것은, 이러한 의료 데이터 침해 사고를 식별하고 통제하는 데 평균 279일이라는 긴 시간이 소요되어 타 산업 대비 피해 규모가 눈덩이처럼 커진다는 점이다.

\bigskip

이처럼 의료계에서 데이터 보안의 중요성이 갈수록 절실해지면서, 근본적인 보안 위협을 줄이기 위한 핵심 원칙으로 '데이터 최소화(Data Minimization)'가 대두되고 있다. 데이터 최소화란 시스템이나 서비스 제공에 반드시 필요한 데이터만을 수집 및 보관하고 그 외의 정보는 철저히 배제하는 것을 의미한다\cite{accessnow2021gdpr}. "수집되지 않은 데이터는 사람들에게 피해를 줄 수 없다"는 원칙 아래, 불필요한 데이터의 수집과 보존을 제한하는 것은 프라이버시 침해를 예방하고 시스템의 보안을 강화하는 가장 사용자 친화적이고 확실한 방법으로 평가받고 있다.

\bigskip

이러한 데이터 최소화 원칙을 시스템적으로 구현하면서도 데이터를 안전하게 유통하기 위해 최근 적극적으로 도입되고 있는 암호학적 기술이 바로 영지식 증명이다. 영지식 증명을 활용하면 개인의 건강 지표나 임상 기록과 같은 민감한 원본 데이터를 외부에 전혀 노출하지 않으면서도, 해당 데이터가 특정 조건을 만족한다는 사실만을 수학적으로 완벽하게 증명할 수 있다. 

\bigskip

\begin{figure}[ht]
    \centering
    \includegraphics[width=1\linewidth]{figure01.png}
    \caption{Gao et al. 외 많은 연구들이 제시한 데이터 증명 시스템의 일반적인 구조}
\end{figure}

하지만 의료 데이터 보안이나 검증 가능한 데이터베이스를 다룬 기존의 많은 연구들은 영지식 증명을 쿼리가 요청될 때마다 조건 일치 여부를 실시간으로 계산하고 버리는 '일시적인 증명 수단'으로만 사용해왔다. 그림 1은 Gao et al.이 제시하는 영지식 증명을 활용한 데이터 거래 시스템\cite{gao2024}의 아키텍쳐를 나타내고 있다. 이는 탈중앙화를 지향하는 데이터 거래 플랫폼 연구들이 대부분 차용하는 일반적인 구조이다. 이 구조에서 데이터 소유자는 데이터를 IPFS(Inter-Planetary File System) 등의 분산 파일 시스템에 저장하고, 해시를 블록체인에 저장하며, 스마트 컨트랙트를 통해 영지식 증명을 수행, 데이터를 재암호화 하여 수요자에게 제공한다. 이 과정에서 실시간 영지식 증명을 통해 소유자는 데이터의 특정 조건 만족 여부를 수요자에게 증명한다. 대부분의 아키텍쳐는 이처럼 데이터 수요가 발생할 때 마다 새로운 증명을 생성하고, 검증 후 폐기하는 구조를 가진다.

\bigskip

이러한 구조는 몇가지 한계를 지닌다. 우선 데이터를 IPFS와 같은 분산 파일 시스템에 암호화 하여 저장한다 하더라도, 결국 소유자가 아닌 제 3자의 시스템에 데이터를 위탁하는 셈이 된다. 이는 이들 연구가 지향하는 데이터 주권의 취지를 약화시키는 요소이다. 또한 스마트 컨트랙트를 통해 데이터 증명과 검증을 수행하는 과정은 소유자와 수요자 간의 1대1 통신으로 이루어진다. 이는 대량의 데이터를 요구하는 기업 및 기관 등에서 활용하는데 있어, 확장성의 문제를 야기한다. 이에 본 연구는 일회성으로 생성되고 버려지는 영지식 증명 데이터들을 한데 모아, 수요자들이 데이터의 존재 여부만 '검색' 할 수 있는, 따라서 오직 소유자의 단말기에만 데이터 원문을 남겨두는 안전한 '증명 검색 시스템(Proof Search System)'을 제안하고자 한다.

\subsection{문제 정의}

\textbf{데이터 프라이버시와 검색의 이중성:} 의료 데이터와 같은 민감한 개인정보를 거래하는 플랫폼이 상용화되기 위해서는, 수요자가 원하는 조건에 부합하는 데이터가 플랫폼 내에 얼마나 존재하는지 확인할 수 있는 '검색' 기능이 필수적이다. 동시에 민감한 정보인 만큼, 검색에 반드시 필요한 속성을 제외한 데이터의 다른 부분들은 철저히 숨겨져야 한다.

\bigskip

그러나 검색과 데이터 프라이버시는 본질적으로 상극의 이중성을 지닌다. 데이터의 유용성은 정보량이 높을수록 증가하지만 그만큼 프라이버시는 감소한다. 마찬가지로 검색의 성립을 위해선 데이터의 일부를 공개를 해야하지만, 그 순간 개인의 프라이버시는 침해될 수 있는 본질적인 모순이 발생한다. 결과적으로, "어떻게 하면 데이터의 원문 그 자체, 혹은 검색에 필수적이지 않은 부분을 은닉하여 데이터 최소화를 실현하면서도, 실시간 처리가 가능한 수준의 검색 효율성을 동시에 달성할 수 있는가?"가 본 연구가 해결하고자 하는 핵심 문제이다.

\subsection{연구 목적}

데이터 거래 플랫폼은 다양한 세부 프로세스들의 집합으로 이뤄진다. 플랫폼 내에서 데이터의 검증, 결제, 전달, 가공, 위탁 등 수많은 시나리오들이 유기적이고 무결하게 이뤄져야 한다. 본 연구는 그러한 시나리오 가운데 가장 첫 단계라 할 수 있는 '데이터 존재 여부의 확인', 즉 검색에 대한 솔루션을 제공한다. 그리고 이를 위해 다음과 같은 핵심 아이디어들을 활용한다.

\bigskip

\textbf{① 원문을 저장하지 않는 증명 검색 시스템:} 일반적인 검색 시스템과는 달리 본 연구는 데이터의 원문을 일절 저장하지 않는다. 대신 원문의 수치가 특정 범위 안에 존재한다는 범위 증명(ZKRP, Zero Knowledge Range Proof) 데이터를 저장하여 검색 대상으로 한다. 의료 영역에서 특정 수치가 일정 범위 안에 존재하는지의 여부는 중요한 진단 지표가 된다. 그리고 이러한 수치와 임계값 사이의 관계는 단순한 디지털 서명 등으로는 증명이 불가능하다. 여기서 영지식 증명을 활용하면 범위의 포함 여부를 확인하면서 동시에 정확한 수치를 공개하지 않는 것이 가능하다. 이를 통해 데이터 수요자는 원하는 데이터의 ‘신뢰 가능한 존재 여부’를 제공받고, 데이터 소유자는 원문의 유출에 대한 우려를 완벽히 제거할 수 있다.

\bigskip

\textbf{② 사전 정의된 참고치 기반의 오프라인 증명 생성:} 실시간 증명 생성의 병목을 해소하기 위해, 본 연구는 의료 데이터 도메인에 특화된 속성을 활용한다. 수치형 의료 데이터는 일반적으로 진단에 활용되는 절대적인 임상 참고치가 존재한다. 예를 들어, 세계보건기구(WHO) 가이드라인\cite{who2011}에 따르면 당화혈색소(HbA1c) 수치는 정상(5.7\% 미만), 당뇨 전 단계(5.7~6.4\%), 당뇨병(6.5\% 이상)으로 구간이 명확히 구분된다. 본 시스템은 이러한 사전 정의된 범위를 바탕으로 데이터 발급 시점에 영지식 증명을 미리 생성하여 서버에 저장한다. 이를 통해 검색 시점에는 실시간 회로 연산 없이 단순 DB 조회만으로 결과를 반환하여 검색 속도를 획기적으로 향상시킨다.

\bigskip

\textbf{③ 단일 그룹 커밋을 통한 복합 쿼리 소유권 증명:} 수요자가 복합 조건(예: 당화혈색소 6.5\% 이상 AND 공복혈당 126 이상)을 쿼리로 하여 검색할 때, 반환된 다수의 증명 데이터들이 모두 동일한 환자(소유자)의 것임을 보장해야 한다. 이는 ‘특정 환자가 어떤 데이터를 갖고 있는가’ 와 같은 환자에 대해 필수적이지 않은 추가 정보를 노출하는 문제를 발생시킨다. 본 연구는 이에 대한 해결책으로 개별 증명의 식별자들을 해싱하여 하나의 그룹으로 묶고(prf\_hash), 이를 소유자의 식별자(DP\_ID) 및 난수(nonce)와 결합하여 단일한 커밋을 생성한다. 이 방식을 통해 악의적인 공격자가 서로 다른 소유자의 데이터를 섞는 조작을 수학적으로 원천 차단하며, 그룹 내 증명 개수와 무관하게 단 1회의 증명 생성만으로 소유권을 입증하는 확장성을 달성한다.

\bigskip

이러한 접근을 통해 본 연구는 데이터 원문을 제 3자인 서버가 보유하지 않는 '완전한 영지식성'을 유지하면서도, 검색 조건에 부합하는 데이터의 존재를 증명할 수 있는 프라이버시와 검색 가능성 사이의 절충안을 제시한다. 이를 바탕으로 누구든 유출 걱정 없이 의료 데이터를 안전하게 판매 및 구매할 수 있는 신뢰도 높은 데이터 유통망 구축의 발판을 마련하고자 한다.